\noindent\textbf{Elección del procedimiento.} La continuidad da $A_1V_1=A_2V_2$ y, por tanto, $V_1=\beta^2V_2$. Bernoulli horizontal sin pérdidas da $\Delta p=\rho(V_2^2-V_1^2)/2=\rho V_2^2(1-\beta^4)/2$. Despejar primero velocidad y después multiplicar por el área de garganta.

\begin{enumerate}

\item El diámetro determina el área de garganta.

\[\fitmath{A_{2}=\frac{\pi\cdot {d_{2}}^{2}}{4}}\]

\[\fitmath{A_{2}=\frac{\pi\cdot {(0.05)}^{2}}{4}\approx 0.001963495\,\mathrm{m^{2}}}\]

\item Tomar la raíz positiva del balance de energía.

\[\fitmath{V_{2}=\sqrt{\frac{2\cdot \mathrm{dp}}{\rho\cdot \left(1-{\beta}^{4}\right)}}}\]

\[\fitmath{V_{2}=\sqrt{\frac{2\cdot (20000)}{(998)\cdot \left(1-{(0.5)}^{4}\right)}}\approx 6.538514\,\mathrm{m/s}}\]

\item Aplicar $Q=A_2V_2$ en la misma sección.

\[\fitmath{Q=A_{2}\cdot V_{2}}\]

\[\fitmath{Q=(0.001963495408)\cdot (6.538514434)\approx 0.01283834\,\mathrm{m^{3}/s}}\]

\end{enumerate}
